\section{Metodología}

El desarrollo del proyecto se basa en la integración de dos marcos metodológicos complementarios: CRISP-DM (Cross Industry Standard Process for Data Mining) \cite{crispdm1999}, ampliamente reconocido en el ámbito del análisis de datos e inteligencia artificial, y la guía TRIPOD (Transparent Reporting of a multivariable Prediction model for Individual Prognosis or Diagnosis) \cite{collins2015}, orientada a garantizar transparencia, reproducibilidad y rigor metodológico en el desarrollo de modelos predictivos clínicos.

El uso conjunto de ambos enfoques permite estructurar el proyecto desde una perspectiva técnico-científica, asegurando no solo la calidad del modelo, sino también su validez clínica y su potencial de integración en entornos hospitalarios.

\subsection{Fases metodológicas}

\textbf{Fase 1: Comprensión del problema y objetivos clínicos}

Corresponde a la formulación inicial del problema, el análisis del contexto clínico y la definición de los objetivos específicos del modelo predictivo. Durante esta fase se identifican los interesados principales (anestesiólogos, investigadores clínicos y especialistas en IA), los criterios de éxito clínico, así como las restricciones éticas y técnicas asociadas al manejo de datos sensibles. Esta fase se desarrolla siguiendo la primera etapa de CRISP-DM Business Understanding y las directrices iniciales de TRIPOD relativas a la definición clara del objetivo de predicción y la población de estudio. Se utiliza Microsoft Word para la documentación formal, Sharepoint para control colaborativo de versiones y Mendeley para la gestión de referencias.

\textbf{Fase 2: Comprensión y preparación de los datos}

En esta etapa se realiza el análisis exploratorio y la preparación del conjunto de datos proporcionado por la Fundación Valle del Lili (FVL), siguiendo las etapas Data Understanding y Data Preparation de CRISP-DM.

Las actividades incluyeron:

\begin{itemize}
    \item Exploración de distribución de variables clínicas y manejo de valores atípicos.
    \item Limpieza, transformación y codificación de variables categóricas y numéricas.
    \item Análisis de correlaciones y selección de variables relevantes con base en criterios clínicos y estadísticos.
\end{itemize}

De acuerdo con TRIPOD, en esta fase también se documentaron las características de la población, los métodos de adquisición de datos, los criterios de exclusión y el tratamiento de valores faltantes, garantizando la trazabilidad de los procedimientos. Las herramientas utilizadas son: Python (con librerías Pandas, NumPy, Matplotlib, Seaborn, entre otras) y Jupyter Notebook para el análisis exploratorio y la limpieza de datos; GitHub para control de versiones.

\textbf{Fase 3: Modelado y experimentación}

Esta fase corresponde a las etapas Modeling y Evaluation de CRISP-DM, y constituye el núcleo del trabajo actual. En ella se compararán distintos algoritmos de aprendizaje supervisado entre ellos Random Forest, Gradient Boosting (LightGBM) con el fin de determinar el modelo de mejor desempeño para la predicción temprana de shock hemorrágico.

Las actividades contempladas son:

\begin{itemize}
    \item División del dataset en conjuntos de entrenamiento, validación y prueba.
    \item Ajuste de hiperparámetros mediante validación cruzada estratificada.
    \item Evaluación del desempeño con métricas clínicas: sensibilidad, especificidad, F1-score, AUC-ROC y calibración, con enfoque en alta sensibilidad
    \item Interpretabilidad del modelo mediante técnicas SHAP para garantizar la transparencia y la validación clínica de los resultados.
\end{itemize}

Desde la perspectiva TRIPOD, esta fase incluye la documentación exhaustiva del proceso de modelado, el método de validación interna y el análisis de robustez y sesgos.

\textbf{Fase 4: Validación clínica preliminar}

En esta etapa se realiza una evaluación preliminar del modelo junto con expertos médicos, utilizando escenarios quirúrgicos históricos simulados. El propósito es verificar la coherencia clínica de las predicciones y la aplicabilidad práctica del modelo en entornos reales.

Esta fase sigue las recomendaciones TRIPOD sobre validación y reporte de resultados, incluyendo:

\begin{itemize}
    \item Identificación de limitaciones y posibles fuentes de sesgo.
    \item Elaboración de tablas y figuras explicativas conforme a los lineamientos de reporte TRIPOD.
\end{itemize}

\textbf{Fase 5: Despliegue y Documentación final}

Se consolida la documentación técnica y científica del proyecto, incluyendo el informe final, las visualizaciones de resultados y los anexos de código reproducible. En este caso, no se integrará el modelo con los sistemas de FVL, pero se busca realiza la etapa Deployment de CRISP-DM con versionamiento de modelos.

\newpage

% -----------------------------------------------------------------------------
% 5. DESARROLLO DE LA PROPUESTA
% -----------------------------------------------------------------------------
